\section{Scope and the exact arithmetic return}
The preceding integral-completion calculation retained the original finite
orders, but produced rational numerical candidates which pass a specified finite
set of local tests while violating one original prime-exponent bound
\cite[\S\,``One unavailable prime occurrence'']{Previous}. We now test a global
observation instead. For rational ES denominators, two integral Newton traces
force all denominators to be integers. In the original raw integer gate source,
one integral trace already forces the correct divisor box, with a necessary
middle-to-exterior reclassification at the distinguished prime. The complete
trace fibres recover the middle states as actual kernel generators.

This is an arithmetic integrality theorem, not an assertion that an integral
trace is attained at every prime. The remaining universal existence question is
stated at the end. In particular, an algebraic-integer root is not silently
identified with a rational integer. We construct an explicit algebraic
counterfamily at every hard prime which retains all integral power traces and
the middle prime-local signature, but has no rational denominator pair.
No historical-priority determination is claimed.

\section{A denominator theorem for finite rational spectra}
Let $x_1,\ldots,x_n\in\Q^\times$, $n\ge2$. Write
\[
 s_k=\sum_{i=1}^n x_i^k,\qquad
 \sum_{i=1}^n\frac1{x_i}=\frac{m}{b}\ne0,
 \qquad b>0,\quad\gcd(m,b)=1.
\]
The numerator $m$ is signed; divisibility of it means divisibility of $|m|$.
The roots are labelled and may repeat.

\begin{theorem}[Common denominator and reciprocal numerator]\label{thm:general}
If $s_1,\ldots,s_{n-1}\in\Z$, the reduced denominators of all $x_i$ are the same
integer $d$. Moreover
\[
 \boxed{\gcd(d,b)=1,\qquad d^n\mid(n-1)!m.}                 \tag{TR1}
\]
If $n$ is odd, $d$ is odd.
\end{theorem}
\begin{proof}
Let $e_j$ be the elementary symmetric functions, with $e_0=1$. The exact Newton
identities are
\[
 j e_j=\sum_{k=1}^j(-1)^{k-1}e_{j-k}s_k.                 \tag{TR2}
\]
They follow by expanding the logarithmic derivative of
$\prod_i(1+x_iT)$ as a formal power series and comparing the coefficient of
$T^{j-1}$. This is the classical identity; compare Euler
\cite[\S\S3,5--7]{Euler}. Induction in TR2 proves
$j!e_j\in\Z$ for $j<n$: multiplying its right side by $(j-1)!$ clears all
previous denominators.

Fix a prime $\ell$ at which some $x_i$ is nonintegral. Let $-a<0$ be their
minimum valuation, and let exactly $r$ roots have this valuation. If $r<n$, the
product of those $r$ roots is the unique least-valuation term in $e_r$, so
$v_\ell(e_r)=-ar$. But
\[
 v_\ell(e_r)\ge-v_\ell(r!)>-r\ge-ar,
\]
a contradiction. The strict factorial bound follows from
$\sum_{k\ge1}\lfloor r/\ell^k\rfloor<r$ (or its finite base-$\ell$ digit
formula). Therefore every root has the same valuation $-a$ at that prime.
This proves equality of all reduced denominators.

Write $x_i=A_i/d$, where each $\gcd(A_i,d)=1$. Since
$e_{n-1}/e_n=m/b$ and $e_n=\prod_iA_i/d^n$, integrality of
$(n-1)!e_{n-1}$ gives at every $\ell^a\parallel d$
\[
 na\le v_\ell((n-1)!)+v_\ell(m)-v_\ell(b).              \tag{TR3}
\]
The factorial valuation is less than $n$. Hence $\ell\mid m$; reducedness then
gives $\ell\nmid b$. Equation TR3 proves TR1. Finally, if $d$ were even and
$n$ odd, all reduced numerators would be odd, and their sum would be odd,
contradicting $s_1\in\Z$.
\end{proof}

\begin{corollary}[Numerator four, any length at least three]\label{cor:four}
Let $n\ge3$, $p\in\Z_{>0}$, and $x_i\in\Q_{>0}$ satisfy
$\sum_i1/x_i=4/p$. Then
\[
 \boxed{x_i\in\Z\ (1\le i\le n)
 \quad\Longleftrightarrow\quad s_1,\ldots,s_{n-1}\in\Z.} \tag{TR4}
\]
For the Erd\H{o}s--Straus equation, only $s_1,s_2$ are required. If $p$ is odd,
these also imply $s_2\equiv s_1^2\pmod8$.
\end{corollary}
\begin{proof}
The reduced reciprocal numerator divides four. An odd prime cannot divide $d$,
since $v_\ell((n-1)!)<n$ in TR1. If $n$ is odd, the final assertion of
Theorem~\ref{thm:general} also excludes two. If $n\ge4$ is even, $n-1$ is odd
and at least three. Counting the binary digits of $n-1$ gives
\[
 v_2((n-1)!)+2=n+1-s_2^{\rm bin}(n-1)\le n-1,
\]
so TR1 again excludes two. Thus $d=1$. The converse is literal. For three
integer denominators, $4e_3=pe_2$ and odd $p$ give $4\mid e_2$; use
$s_2=s_1^2-2e_2$.
\end{proof}
The restriction $n\ge3$ is necessary: $(x_1,x_2)=(p/2,p/2)$ at odd $p$ has
integer first trace and reciprocal sum $4/p$, but is not integral.

\subsection{The exact possible denominators at length three}
\begin{theorem}[Sharp three-root denominator classification]\label{thm:three}
Under Theorem~\ref{thm:general} with $n=3$,
\[
 \boxed{d\text{ is odd},\quad v_3(d)\le1,\quad
 \ell\mid d,\ \ell\ne3\Longrightarrow\ell\equiv1\pmod3,
 \qquad d^3\mid m.}                                   \tag{TR5}
\]
Conversely every $d$ with these prime conditions is realized by three positive
rationals with common reduced denominator $d$, integral $s_1,s_2$, and
$v_\ell(m)=3v_\ell(d)$ at every $\ell\mid d$.
\end{theorem}
\begin{proof}
The oddness follows above, and TR1 then gives $d^3\mid m$. At a denominator
prime, scale the three roots to local units $A,B,C$. Their first and second
symmetric functions vanish modulo $\ell$, so their residues are the roots of
$T^3-ABC$. If $\ell\ne3$, the polynomial has distinct nonzero roots. Ratios of
these roots supply a nontrivial cube root of one, giving $\ell\equiv1\pmod3$.
For $\ell=3$ their residues coincide. If $a=v_3(d)\ge2$, the exact symmetric
identities imply
\[
 A^2+AB+B^2=3^a s_1(A+B)-3^{2a}e_2.
\]
The right side has valuation at least two, while
$4(A^2+AB+B^2)=(2A+B)^2+3B^2$ has valuation exactly one. This is impossible.
Other denominator primes are units in this displayed local equation.

The denominator $d=1$ is realized by $(1,2,3)$. For $d>1$, work independently at each $\ell^a\parallel d$ modulo
$\ell^{2a+1}$, taking the first numerator to be one. At $\ell=3$, $a=1$, take
$(A,B,C)=(1,1,4)\pmod{27}$; the sum is divisible by three and the pair sum has
valuation exactly two. At any other denominator prime choose a nontrivial cube
root $b$ modulo $\ell^{2a}$. It is obtained from a root of $b^2+b+1$ modulo
$\ell$ by the explicit unit-derivative lift
\[
 b'=b+\ell^k t,\qquad
 t\equiv-\frac{b^2+b+1}{\ell^k}(2b+1)^{-1}\pmod\ell.
\]
Set $c=-1-b+\ell^{2a}t$. Its pair sum is
$-(b^2+b+1)+(1+b)\ell^{2a}t$. Because $1+b$ is a unit, precisely one value
of $t\pmod\ell$ makes its valuation exceed $2a$; choose any other value.
CRT now gives integer $B,C$ with these separate labelled residues. Add
respectively one and two multiples of
$L=\prod_{\ell^a\parallel d}\ell^{2a+1}$ to make $1<B<C$.
All three numerators are coprime to $d$; their sum is divisible by $d$ and their
pair sum has exact valuation $2a$ at each denominator prime. Dividing by $d$
gives integral first two traces. The reciprocal numerator has exact valuation
$3a$ because the product of the numerators is a unit there.
\end{proof}
Examples retaining the exact reciprocal numerators are
\[
 \begin{array}{c|c|c}
 (x_1,x_2,x_3)&(s_1,s_2)&\sum_i1/x_i\\\hline
 (1/3,1/3,4/3)&(2,2)&27/4\\
 (1/7,18/7,30/7)&(7,25)&343/45\\
 (1/13,22/13,146/13)&(13,129)&10985/803.
 \end{array}                                           \tag{TR6}
\]
The cube restriction concerns the reduced numerator, not the number of primes
one has elected to check.

\section{One global trace on an unfiltered integer gate source}
Now fix prime $p\equiv1\pmod4$ and a complete original first-half shell
\[
 p/4<a<p/2,\qquad R=4a-p.
\]
Then $R\ge3$, $R\equiv3\pmod4$, and $\gcd(pa,R)=1$.
Before imposing any divisor budget, take an arbitrary positive integer $U$ with
$R\mid a+U$. Define
\[
 \boxed{x=a,\qquad y=\frac{pa(U+a)}{RU},\qquad
 z=\frac{p(a+U)}R,\qquad S=p+x+y+z.}                    \tag{TR7}
\]
These are positive rational denominators, $z$ is an integer, and direct
multiplication gives $1/x+1/y+1/z=4/p$.

\begin{theorem}[Global trace recovers the original budget and tag]\label{thm:raw}
The exact reduced denominator is
\[
 \boxed{\den S=\den y=\frac{U}{\gcd(U,pa^2)}.}           \tag{TR8}
\]
Consequently $S\in\Z$ if and only if $U\mid pa^2$. Its integral returns are
exactly the disjoint tagged cases
\[
 \begin{array}{c|c|c}
 v_p(U)&\text{original divisor}&\text{channel}\\\hline
 0&u=U\mid a^2&M\\
 1&u=U/p\mid a^2&E.
 \end{array}                                           \tag{TR9}
\]
No higher $p$-valuation gives an integral return. On an integral trace the tag
can also be read without factoring $U$:
\[
 \boxed{M\iff S\equiv a\pmod p,\qquad E\iff S\not\equiv a\pmod p.}
                                                               \tag{TR10}
\]
For the raw E gate $R\mid4U+1$, the analogous rational return
\[
 \left(a,\frac{paU+a^2}{RU},\frac{pa+p^2U}{R}\right)
\]
has $\den S=U/\gcd(U,a^2)$; its trace is integral exactly on the original E box.
\end{theorem}
\begin{proof}
The gate makes $a+U$ divisible by $R$, and $\gcd(U,R)=1$. If
$n=pa(U+a)/R\in\Z$, then, modulo $U$,
$Rn\equiv pa^2$. Hence $\gcd(n,U)=\gcd(pa^2,U)$, proving TR8. For integral
returns $a<p$ makes the $p$-part of $U$ zero or one. At valuation zero, the
original M gate is unchanged. At valuation one write $U=pu$. Since
$p\equiv4a\pmod R$, the gate $a+pu\equiv0$ is exactly $4u+1\equiv0$; the two
expressions in TR7 are the original E denominators at $u$. They are not a
canonical M state with an impermissible extra occurrence of $p$.

For completeness, the original normalization is
\[
 g=\gcd(a,u),\quad h=g^2/u,\quad r=u/g,\quad s=a/g.
\]
Prime valuations give positive integral $h,r,s$, $a=hrs$, $u=hr^2$ and
$\gcd(r,s)=1$. At M set $\lambda=(r+s)/R$ and return
$(a,phs\lambda,phr\lambda)$. At E set $\kappa=(pr+s)/R$ and return
$(a,hs\kappa,phr\kappa)$. The gates prove their integrality. The ordered
inverses are $pa^2/(Ry-pa)$ and $a^2/(Ry-pa)$ respectively. These Type-I/II
identities are inherited \cite[\S2, Propositions 2.2 and 2.6]{ET}.
At M both tail denominators are multiples of $p$, so TR10 holds. At E the first
tail denominator is a $p$-unit, proving the other alternative. The raw E
calculation is identical: its third denominator is integral by multiplying the
E gate by $pa$, and reduction of the second numerator modulo $U$ gives $a^2$.
\end{proof}

For a nonintegral raw word, TR8 gives the exact integer-spacing estimate
\[
 \boxed{\operatorname{dist}(S,\Z)\ge\frac1{\den S}\ge\frac1U.}  \tag{TR11}
\]
Strict proximity $|S-N|<1/U$ to an integer forces equality. No estimate proving
that proximity for some raw word has been obtained. The strict sign is necessary:
at $p=13,a=4,R=3,U=5$, $S=436/5$ and the distance is $1/5$.

At $p=35281,a=8821,U=202883$, the preceding foreign-prime candidates have exact
literal traces
\[
 S_E=1936132532779018/23,\qquad S_M=59754044938/23.        \tag{TR12}
\]
Thus the global trace rejects their one unavailable prime occurrence. This is a
stronger \emph{test} than the earlier fixed finite-place predicate, not a proof
that the test succeeds at every prime.

\section{The trace fibres and their integral coefficient kernel}
Let
\[
 \Omega_{p,a}=\{1\le U\le pa^2:R\mid a+U\},
 \quad \Theta_{p,a}=\{S(U):U\in\Omega_{p,a}\},
\]
and define the actual free-abelian pushforward
\[
 \Phi:\Z^{\Omega_{p,a}}\longrightarrow\Z^{\Theta_{p,a}},
 \qquad e_U\longmapsto e_{S(U)}.                         \tag{TR13}
\]
The raw source is finite without assuming any original hit. Its explicit size is
$1+\lfloor(pa^2-U_0)/R\rfloor$, where $U_0=(-a)\bmod R\in[1,R-1]$.

\begin{theorem}[Complete trace fibres]\label{thm:fibre}
Two distinct raw words have the same trace exactly when $UV=a^2$. Every such
pair consists of the two original middle orientations. Every exterior integral
trace has one raw preimage; every middle integral trace has two; every
nonintegral trace has one. If $E_a,M_a$ are the original labelled counts, with
both middle orientations retained, then
\[
 \boxed{\rank\ker\Phi=M_a/2,\quad
 |\Theta_{p,a}\cap\Z|=E_a+M_a/2,\quad
 E_a+M_a=|\Theta_{p,a}\cap\Z|+\rank\ker\Phi.}            \tag{TR14}
\]
The kernel has the explicit basis $e_U-e_{a^2/U}$ for the middle words $U<a$.
\end{theorem}
\begin{proof}
From TR7,
\[
 S(U)-S(V)=\frac pR(U-V)\left(1-\frac{a^2}{UV}\right).  \tag{TR15}
\]
If both raw words are integers and $UV=a^2$, they are original square divisors.
The gate at $U$ implies the gate at $V$ since $U\equiv-a\pmod R$ and is a unit.
Thus both are M. There is no fixed point: $U=a$ would make $R\mid2a$, impossible.
Conversely every original M complement is present in $\Omega_{p,a}$. An E word
has $v_p(U)=1$; its complementary rational number $a^2/U$ is not an integer,
so its fibre is a singleton. These are the entire fibres, proving TR14 and
the stated integral basis.
\end{proof}

The labelled inverse from a trace value $S$ is the positive-root equation
\[
 U^2-\left(\frac R p(S-p-a)-2a\right)U+a^2=0,            \tag{TR16}
\]
followed by the original integer, range and gate tests. Keeping the source word
$U$ recovers its factor exponents and tag through TR9. Equal trace values at
\emph{different} distinguished denominators $a$ are not merged by TR13.

For arbitrary retained positive source weights $w_U$, use
$\|x\|^2=\sum_Uw_U|x_U|^2$. The exact quotient norm and minimum-norm section are
\[
 \|z\|_{\rm quot}^2=\sum_S\frac{|z_S|^2}{\sum_{S(U)=S}w_U^{-1}},
 \qquad x_U=\frac{z_{S(U)}w_U^{-1}}{\sum_{S(V)=S(U)}w_V^{-1}}.
\]
Cauchy--Schwarz applied on each actual fibre proves the lower bound, and
substitution proves equality for the displayed section. In the counting norm
$w_U=1$, it places $z_S/m_S$ on each of the $m_S$ original words, with squared
quotient norm $\sum_S|z_S|^2/m_S$. At a two-point fibre this section is not
integral when $z_S$ is odd. Choosing one original representative instead gives
an integral section with a different norm. No section discards the kernel.
These are norms of the specified finite word source, not a replacement of the
weighted analytic conductor. Reduction
modulo two sends the original middle mass two to zero but retains its two words
and the kernel generator. It is not an absent source.

For example, at $p=1009,a=253$, the middle pair $U=11,5819$ has
$S=2132270$. The exterior word $U=11099=p\cdot11$ has $S=3906350$ and returns
\[
 (253,87032,3818056).
\]
It is integral but is not a canonical M state at the uncorrected word $11099$.

\section{An integral first trace alone retains a square-part obstruction}
Consider arbitrary positive rational tails $y,z$ at the original integer $a$:
\[
 \frac1y+\frac1z=\frac R{pa},\qquad N=y+z\in\Z.
\]
This section does not assume the raw full gate of TR7.

\begin{theorem}[Exact remaining denominator]\label{thm:tail}
The reduced denominators of $y,z$ coincide at $d$, and
\[
 \boxed{d^2=\frac R{\gcd(R,N)},\qquad
 \den(a^2+y^2+z^2)=d^2.}                                \tag{TR17}
\]
In particular squarefree $R$ makes an integral first trace sufficient.
Every such pair has the complete integer-factor presentation
\[
 b c=(pa)^2,\quad b,c>0,\quad
 y=(pa+b)/R,\quad z=(pa+c)/R,
\]
with
\[
 \boxed{K(R)\mid pa+b,\qquad
 K(R)=\prod_{\ell^e\parallel R}\ell^{\lceil e/2\rceil}.} \tag{TR18}
\]
The pair is integral precisely when $R\mid pa+b$; in the relaxed pair
$d=R/\gcd(R,pa+b)$.
\end{theorem}
\begin{proof}
Since $y+z$ is integral, write $y=Y/d,z=Z/d$ with both numerators coprime to $d$.
The reduced denominator of $yz$ is $d^2$. But $yz=Npa/R$ has denominator
$R/\gcd(R,N)$, proving the first equality. The second follows from
$a^2+y^2+z^2=a^2+N^2-2Npa/R$ and odd $R$, coprime to $pa$.
Thus $d\mid R$, and $b=Ry-pa,c=Rz-pa$ are positive integers; the reciprocal
identity gives $bc=(pa)^2$. They are units modulo $R$. Conversely every such
factor pair gives positive rational tails and
\[
 y+z=\frac{2pa+b+c}{R}=\frac{(pa+b)^2}{Rb}.
\]
Its integrality is $R\mid(pa+b)^2$, exactly TR18. The remaining assertions
are the literal denominators of $(pa+b)/R$.
\end{proof}

The failed shell-change attempt is completely concrete at the hard prime $1009$:
\[
 \boxed{a=286,\ R=135,\ K(R)=45,\ b=11,\quad
 (x,y,z)=(286,6413/3,168238642/3).}                      \tag{TR19}
\]
It satisfies ES, has first denominator trace $56081971\in\Z$, and second trace
$28304240703866897/9$. Its missing denominator is exactly three. Removing its
square factor from $R$ gives $R'=15,a'=256$. The complete original divisors are
$2^f$, $0\le f\le16$, whose residues modulo fifteen are $1,2,4,8$.
Neither original target, $14$ for M or $11$ for E, is present. Thus stripping
that square does not transport the original factor box to an occupied shell.
The smaller control $p=193,a=55,R=27$ similarly gives
$(55,1180/3,2505140/3)$ with integral first trace, while the compressed shell
$R'=3,a'=49$ has both gates empty. These are failures of the specified map,
not ES counterexamples.

\section{The literal coefficient observation and a rationality boundary}
For the original literal roots $(p,x,y,z)$, keep the exact normalization
\[
 \mathcal H(T)=AT^4+T^3+BT^2+CT+D
              =-S^{-1}(T-p)(T-x)(T-y)(T-z).
\]
These are the preceding integral-completion coordinates, not the reciprocal
cubic and not the analytic conductor moments \cite[IC4--IC5]{Previous}.
The first two traces of multiplication by $T$ on the rank-four root algebra are
\[
 \boxed{\mathsf N_1=-A^{-1},\qquad
 \mathsf N_2=A^{-2}-2B/A.}                              \tag{TR20}
\]
They follow from TR2. For rational ES roots, Corollary~\ref{cor:four} says that
integrality of these two traces is equivalent to the original integer
solution, since subtracting $p,p^2$ gives the denominator traces. On the rank-eight
completion obtained by adjoining $\eta^2=\mathcal H'(T)$, the traces of these
multiplication operators are \emph{twice} TR20. The two copies and their original
basis are retained; a weighted analytic shift-conductor moment is a different
observation and has not been substituted for either one.

The two traces also determine the whole unordered denominator polynomial, not
merely an integrality predicate. With $s_1=\mathsf N_1-p$ and
$s_2=\mathsf N_2-p^2$, it is exactly
\[
 \boxed{g_{p,s_1,s_2}(T)=T^3-s_1T^2+
 \frac{s_1^2-s_2}{2}T-\frac{p(s_1^2-s_2)}8.}             \tag{TR20a}
\]
The reciprocal equation supplies its constant term. Every integral trace pair
whose polynomial splits into positive rational roots therefore gives integer
roots by Corollary~\ref{cor:four}; in particular $s_1^2-s_2$ is divisible by
eight at odd $p$. On the original canonical domain, choose the unique root
$a\in(p/4,p/2)$. Its remaining $p$-divisibility pattern selects E (one tail
multiple of $p$) or M (both tails multiples of $p$), with the M orientation bit
retained. The original ordered divisor inverse then recovers the entire box
marking. The fibres of the map from original states to
$(p,\mathsf N_1,\mathsf N_2)$ have size one at E and two at M. These claims use
only the original distinct-denominator domain; arbitrary input trace pairs
still require positivity, rational splitting and the original range test.

\begin{theorem}[Algebraic targets with every integer trace at every hard prime]
\label{thm:algebraic}
For every prime $p\ge1009$, $p\equiv1\pmod{24}$, there are explicitly constructed
positive, distinct algebraic-integer denominators $(a,pb,pc)$ satisfying ES,
with $a=(p+3)/4$, for which $b,c$ are irrational over $\Q$. All power traces
are integers. Their literal quartic splits locally over $\Q_p$ and has the
same middle root-cluster valuations and negative pair character as an original
M state. It also satisfies the inherited scalar height and separation bounds.
This family is outside the rational-root domain of Corollary~\ref{cor:four}.
\end{theorem}
\begin{proof}
There are exactly $(p-1)/4$ residues $k_0\pmod p$ such that
\[
 \left(\frac{k_0}{p}\right)=
 \left(\frac{3k_0-1}{p}\right)=-1.                     \tag{TR21}
\]
To count them set $s=3k_0$, since $3$ is a square. The sum
$\sum_s\chi(s(s-1))=-1$ follows by counting
$(2s-1)^2-(2y)^2=1$: there are $p-1$ nonzero choices of its first factor,
each uniquely determining its second. Expand
$(1-\chi(s))(1-\chi(s-1))/4$. The values at zero and one are zero because
$-1$ is a square, so the count is exactly $(p-1)/4$.

Let $k_0$ be the least such residue, and set
\[
 m=(p-1)/6=(2a-2)/3,\quad
 K_0=\lfloor(m+1)^2/2\rfloor+1,\quad
 k=K_0+((k_0-K_0)\bmod p).
\]
Let $b,c$ be the two roots of
\[
 T^2-3kT+ak,
 \qquad \Delta=9k^2-4ak.                               \tag{TR22}
\]
The exact inequalities
\[
 \boxed{(3k-m-1)^2<\Delta<(3k-m)^2}                    \tag{TR23}
\]
follow because the first difference is $2k-(m+1)^2>0$ and the second is
$4k+m^2>0$. The bounding integers are positive. Thus $\Delta$ is not a rational
square; $b,c$ are irrational algebraic integers. They are positive,
$b>a/3$ and $c<3k$, from $b=2ak/(3k+\sqrt\Delta)$.
Their reciprocal identity is
\[
 \frac1a+\frac1{pb}+\frac1{pc}
 =\frac1a+\frac3{pa}=\frac4p.
\]
Their monic cubic is
\[
 T^3-(a+3pk)T^2+pa k(p+3)T-p^2a^2k.                   \tag{TR24}
\]
All power traces are integers by TR2 and its integral recurrence beyond degree
three. This does not make its quadratic factor split over $\Q$.
There is an equally literal algebraic version of the raw source:
\[
 U_\pm=\frac{9k-2a\pm3\sqrt\Delta}{2},\quad
 U_+U_-=a^2,\quad a+U_\pm=3b\ \text{or}\ 3c.            \tag{TR25}
\]
Both $U_\pm$ are positive irrational algebraic integers, since they are the roots
of $U^2-(9k-2a)U+a^2$. They are genuine divisors of $a^2$ in that quadratic
integer ring and pass its sum gate. Also
$U_\pm<9k<p a^2$. Substitution in TR7 returns the two tail orientations.
This is not the source $\Z_{>0}$: extending its scalar ring has supplied two
new words. The field and the conjugation label cannot be discarded.

Modulo $p$, $a=3/4$, and
$\Delta=3k(3k-1)$ is a nonzero square by TR21. Unit-derivative Hensel lifting
therefore gives both roots in $\Z_p$. Their product has character
$\chi(ak)=-1$. They are distinct units, and neither equals one modulo $p$:
that would force $k=4/9$, a square. Thus the literal roots $p,pb,pc$ have all
pairwise difference valuations one, and the singleton $a$ is separated by
units. The derivative valuations are $(2,0,2,2)$. With
$S=p+a+3pk$, the normalized coefficient valuations are $(0,1,2,3)$, by direct
expansion. These are the original M cluster conditions from
\cite[IC7, IC14]{Previous}; that calculation does not force global rationality.
For the lift itself, if $z^2=\Delta\pmod{p^n}$, set
$z'=z+p^nt$ with $t=-(z^2-\Delta)/p^n\,(2z)^{-1}\pmod p$; both compositions
are checked directly. Standard background is \cite[Chapter 7]{Milne}.

Finally
\[
 k\le(p+5)^2/72+p<p^2/60\quad(p\ge1009).
\]
The last comparison is $p^2-410p-125>0$. Put
$s_p=(p^2+3p)/4$ and $B_p=p+s_p(s_p+1)$. Then
\[
 pc<3pk<p^3/20<s_p(s_p+1)/3,\qquad S<B_p.
\]
In increasing order the literal roots are $a,p,pb,pc$, with all adjacent gaps
larger than one. Their Vandermonde product is consequently larger than twelve,
and the product of the three differences from any one root is at least two.
The same scalar inequalities as the preceding integer-source bounds follow:
\[
 \Disc(\mathcal H)>144/B_p^6,\qquad
 2/S\le|\mathcal H'(\ell_i)|<S^2.
\]
No integral-point conclusion is drawn from them.
\end{proof}

At $p=1009$ the construction gives
\[
 a=253,\quad k_0=13,\quad k=15148,\quad
 \Delta=2049827360,\quad
 45275^2<\Delta<45276^2.
\]
The quadratic is $T^2-45444T+3832444$. These explicit irrational denominator
roots have all integer traces and the required local split pattern. They are
not original divisor states. In particular, replacing a rational-root test by
an arbitrary finite list of symmetric integer traces is unsound: this example
has \emph{every} symmetric power trace integral.

\section{What has and has not closed}
The global rational-source integrality return is complete. It may be tested by
two traces on arbitrary rational ES denominators, or one trace on TR7 with the
full integer gate and its distinguished prime bit. The exact coefficient kernel
in TR13 also recovers all oriented M states. These maps preserve all original
fibre multiplicities.

The remaining assertion is still that every prescribed hard prime produces
at least one integral trace in that raw source, equivalently an original
opposite divisor pair. No lower bound on $|\Theta_{p,a}\cap\Z|$ or on
$\rank\ker\Phi$ has been proved that guarantees this for some shell. TR17--19
show why stripping a square from a residual is not an available replacement.
Theorem~\ref{thm:algebraic} shows why the algebraic completion cannot supply the
missing rational-root condition by itself. Turn 7 is not marked complete and
no witness-guaranteed Turn 8 procedure is asserted.

The executable bounds, all source tables, rejected false implications, and
normal-mode, optimized-mode and separate-checker outcomes are recorded separately. Finite scans
check the stated formulas and fibres; they do not establish universal occupancy.
